\chapter{Open bridges and the final target}

\begin{definition}[Computational inputs for the current phase]
\label{bp:computational-inputs}
The fifteen computational obligation families listed in the project report
are assumed at the user's request for the current conditional development.
They include specified diagram-2, diagram-3 and diagram-6 classifications,
jet identities, nonvanishing and isolation assertions. The public v10 repository has now been located and all declared Python
checks replay successfully. Except for the P66 subset, these computations
have not been replayed in Lean and are not global Lean axioms.
\end{definition}

\begin{theorem}[Positive algebraic curve bridge: open]
\label{bp:curve-bridge}
\uses{bp:configuration, bp:fourproduct-fibre, bp:virial, bp:tangent-energy, bp:analytic-energy}
\notready
The outstanding construction must take an infinite positive normalized locus
to a suitable irreducible complex curve, prove constancy of the potential on
that curve, and retain a positive nonzero value. The checked identity $U=I$
and the newly checked differentiable-path constancy theorem do not yet
establish global constancy on every irreducible algebraic curve. Nash-arc,
Zariski-closure and connectedness arguments remain geometric obligations.
\end{theorem}

\begin{theorem}[Completion and order bridge: open]
\label{bp:completion-bridge}
\uses{bp:curve-bridge, bp:fourproduct-fibre}
\notready
Construct one completed normalized curve and its consistent meromorphic order
functions. Prove finite support, degree zero, and that a nonconstant
four-product has a pole. Relate the local Puiseux parameters to those orders;
when using finite covers, prove the positive ramification factor that permits
orders from different local charts to be compared on the same original curve.
\end{theorem}

\begin{theorem}[Local diagram bridge: open]
\label{bp:local-bridge}
\uses{bp:computational-inputs, bp:completion-bridge, bp:order-arithmetic, bp:pair-inertia, bp:d7-inertia}
\notready
Formalize the external sixteen-diagram classification and the required local
alternatives, including boundedness restricted to diagrams 10--16 and the
separate diagram-9 estimate. Show that every four-product pole gives a positive
total distance order. The accepted computational inputs do not replace the
unwritten geometric links between chart solutions, germs and global orders.
In particular, the checked diagram-7 polynomial decomposition still needs
the local estimates that make its terms tend to zero in the required regime.
\end{theorem}

\begin{theorem}[Five-body finiteness: final target, open]
\label{bp:main-target}
\uses{bp:conditional-global, bp:completion-bridge, bp:local-bridge, bp:upstream-adapter}
\lean{Smale5.FiveBodyFiniteness}
\notready
For each positive mass vector, the set of positive normalized configurations
is finite. The Lean name here denotes a proposition definition, not a proved
theorem. Even under the accepted computational inputs, the curve, completion
and local-diagram bridges above remain to be proved. Equivalence with the
original physical problem modulo translations, rotations and positive scaling
is partly supplied by the checked upstream adapter. The independent complete
upstream proof is a different route; its full local certificate rebuild and
final integration check remain separately tracked.
\end{theorem}


